\documentclass{article}
\usepackage{amsmath, amssymb}
\usepackage{geometry}

\geometry{a4paper, margin=1in}

\title{The DD(DIV-DEC and DEC-DIV) Operator}
\author{Govis}

\begin{document}

\maketitle

\begin{center}
We can choose a function operator $ dd(x, y) $ . The set $ Q $ initially contains $ 1 $ . \\
If $ Q $ contains $ p $ and $ q $ ,then $ dd(p,q) $ is added to $ Q $ .
\end{center}

\section*{DIV-DEC}

\subsection*{operators}
\begin{align*}
dd(x, y) &= \frac{x}{y} - 1 \\
dec(x) &= x - 1 = dd(x, 1) \\
neg(x) &= -x = dd(dec(x), -1) \\
inc(x) &= x + 1 = neg(dd(x, -1)) \\
div(x, y) &= x / y = inc(dd(x, y)) \\
rec(x) &= \frac{1}{x} = div(1, x) \\
mul(x, y) &= x * y = div(x, rec(y)) \\
sub(x, y) &= x - y = mul(dd(x, y), y) \\
add(x, y) &= x + y = sub(x, neg(y))
\end{align*}


\subsection*{constants}
\begin{align*}
&1 \\
&-1 = dec(dec(1))
\end{align*}

\section*{DEC-DIV}

\subsection*{operators}
\begin{align*}
dd(x, y) &= \frac{x - 1}{y} \\
dec(x) &= x - 1 = dd(x, 1) \\
neg(x) &= -x = dec(dd(x, -1)) \\
inc(x) &= x + 1 = dd(neg(x), -1) \\
div(x, y) &= x / y = dd(inc(x), y) \\
rec(x) &= \frac{1}{x} = div(1, x) \\
mul(x, y) &= x * y = div(x, rec(y)) \\
sub(x, y) &= x - y = dd(div(x, y), rec(y)) \\
add(x, y) &= x + y = sub(x, neg(y))
\end{align*}

\subsection*{constants}
\begin{align*}
&1 \\
&-1 = dec(dec(1))
\end{align*}

\section*{Question}
Can you prove that Q is the set of all rational numbers?

\end{document}